\section{AMC 12 12B 2007}
        \begin{enumerate}
        	\item (\textit{AMC 12 12B 2007}) Isabella's house has 3 bedrooms. Each bedroom is 12 feet long, 10 feet wide, and 8 feet high. Isabella must paint the walls of all the bedrooms. Doorways and windows, which will not be painted, occupy 60 square feet in each bedroom. How many square feet of walls must be painted? 


 $\mathrm{(A) \ } 678\qquad \mathrm{(B) \ } 768\qquad \mathrm{(C) \ } 786\qquad \mathrm{(D) \ } 867\qquad \mathrm{(E) \ } 876$

 \textbf{Answer: }E



\textbf{Solution}

There are four walls in each bedroom (she can't paint floors or ceilings). Therefore, we calculate the number of square feet of walls there is in one bedroom:
 \[2\cdot(12\cdot8+10\cdot8)-60=2\cdot176-60=292\]We have three bedrooms, so she must paint $292\cdot3=\boxed{\textbf{(E) }876}$ square feet of walls.

 


	\item (\textit{AMC 12 12B 2007}) A college student drove his compact car 120 miles home for the weekend and averaged 30 miles per gallon. On the return trip the student drove his parents' SUV and averaged only 20 miles per gallon. What was the average gas mileage, in miles per gallon, for the round trip?


 $\mathrm {(A)} 22\qquad \mathrm {(B)} 24\qquad \mathrm {(C)} 25\qquad \mathrm {(D)} 26\qquad \mathrm {(E)} 28$

 \textbf{Answer: }B



\textbf{Solution 1}

The trip was $240$ miles long and took $\dfrac{120}{30}+\dfrac{120}{20}=4+6=10$ gallons. Therefore, the average mileage was $\dfrac{240}{10}= \boxed{\textbf{(B) }24}$




\textbf{Solution 2}

Alternatively, we can use the harmonic mean to get $\frac{2}{\frac{1}{20} + \frac{1}{30}} = \frac{2}{\frac{1}{12}} = \boxed{\textbf{(B) }24}$



	\item (\textit{AMC 12 12B 2007}) The point $O$ is the center of the circle circumscribed about triangle $ABC$, with $\angle BOC = 120^{\circ}$ and $\angle AOB = 140^{\circ}$, as shown. What is the degree measure of $\angle ABC$?


 
                    \begin{figure}[H]
                        \centering
                        \includegraphics[width=0.5\textwidth]{Problem/AMC_12_12B_2007_3_Problem_Figure_1.png}
                    \end{figure}
                    
$\mathrm {(A)} 35\qquad \mathrm {(B)} 40\qquad \mathrm {(C)} 45\qquad \mathrm {(D)} 50\qquad \mathrm {(E)} 60$

 \textbf{Answer: }D



\textbf{Solution}

Since triangles $ABO$ and $BOC$ are isosceles, $\angle ABO=20^o$ and $\angle OBC=30^o$. Therefore, $\angle ABC=50^o$, or $\mathrm{(D)}$.

 \textbf{Alternative Solution} $\angle AOC = 100^{\circ} \implies \angle ABC =\frac{\angle AOC}{2} =50^{ \circ},$ or $\mathrm{(D)}.$



	\item (\textit{AMC 12 12B 2007}) At Frank's Fruit Market, 3 bananas cost as much as 2 apples, and 6 apples cost as much as 4 oranges. How many oranges cost as much as 18 bananas? 


 $\mathrm {(A)} 6\qquad \mathrm {(B)} 8\qquad \mathrm {(C)} 9\qquad \mathrm {(D)} 12\qquad \mathrm {(E)} 18$

 \textbf{Answer: }B



\textbf{Solution}

$3$ bananas cost as much as $2$ apples, so $18$ bananas cost as much as $12$ apples. Since $6$ apples cost as much as $4$ oranges, $12$ apples cost as much as $8$ oranges, giving $(B)$.

 


	\item (\textit{AMC 12 12B 2007}) The 2007 AMC 12 contests will be scored by awarding 6 points for each correct response, 0 points for each incorrect response, and 1.5 points for each problem left unanswered. After looking over the 25 problems, Sarah has decided to attempt the first 22 and leave the last 3 unanswered. How many of the first 22 problems must she solve correctly in order to score at least 100 points? 


 $\mathrm {(A)} 13\qquad \mathrm {(B)} 14\qquad \mathrm {(C)} 15\qquad \mathrm {(D)} 16\qquad \mathrm {(E)} 17$

 \textbf{Answer: }D



\textbf{Solution}

She must get at least $100 - 4.5 = 95.5$ points, and that can only be possible by answering at least $\lceil \frac{95.5}{6}\rceil = 16 \Rightarrow \mathrm {(D)}$ questions correctly.

 


	\item (\textit{AMC 12 12B 2007}) Triangle $ABC$ has side lengths $AB = 5$, $BC = 6$, and $AC = 7$. Two bugs start simultaneously from $A$ and crawl along the sides of the triangle in opposite directions at the same speed. They meet at point $D$. What is $BD$? 


 $\mathrm {(A)} 1\qquad \mathrm {(B)} 2\qquad \mathrm {(C)} 3\qquad \mathrm {(D)} 4\qquad \mathrm {(E)} 5$

 \textbf{Answer: }D



\textbf{Solution}


                    \begin{figure}[H]
                        \centering
                        \includegraphics[width=0.5\textwidth]{Solution/AMC_12_12B_2007_6_Solution_Figure_1.png}
                    \end{figure}
                    

 One bug goes to $B$. The path that he takes is $\dfrac{5+6+7}{2}=9$ units long. The length of $BD$ is $9-AB=9-5=4 \Rightarrow \mathrm {(D)}$



	\item (\textit{AMC 12 12B 2007}) All sides of the convex pentagon $ABCDE$ are of equal length, and $\angle A = \angle B = 90^{\circ}$. What is the degree measure of $\angle E$? 


 $\mathrm {(A)} 90\qquad \mathrm {(B)} 108\qquad \mathrm {(C)} 120\qquad \mathrm {(D)} 144\qquad \mathrm {(E)} 150$

 \textbf{Answer: }E



\textbf{Solution}


                    \begin{figure}[H]
                        \centering
                        \includegraphics[width=0.5\textwidth]{Solution/AMC_12_12B_2007_7_Solution_Figure_1.png}
                    \end{figure}
                    

 Since $A$ and $B$ are \href{https://artofproblemsolving.com/wiki/index.php?title=Right\_angle}{right angles}, and $AE$ equals $BC$, and $AECB$ is a \href{https://artofproblemsolving.com/wiki/index.php?title=Square}{square}. Since the length of $ED$ and $CD$ are also 5, triangle $CDE$ is \href{https://artofproblemsolving.com/wiki/index.php?title=Equilateral\_triangle}{equilateral}. Angle $E$ is therefore $90+60=150 \Rightarrow \mathrm {(E)}$



	\item (\textit{AMC 12 12B 2007}) Tom's age is $T$ years, which is also the sum of the ages of his three children. His age $N$ years ago was twice the sum of their ages then. What is $T/N$ ? 


 $\mathrm {(A)} 2\qquad \mathrm {(B)} 3\qquad \mathrm {(C)} 4\qquad \mathrm {(D)} 5\qquad \mathrm {(E)} 6$

 \textbf{Answer: }D



\textbf{Solution}

Tom's age $N$ years ago was $T-N$. The sum of the ages of his three children $N$ years ago was $T-3N,$ since there are three children. If his age $N$ years ago was twice the sum of the children's ages then,
 \begin{align*}T-N&=2(T-3N)\\ T-N&=2T-6N\\ T&=5N\\ T/N&=\boxed{\mathrm{(D) \ } 5}\end{align*}
Note that actual values were not found.

 


	\item (\textit{AMC 12 12B 2007}) A function $f$ has the property that $f(3x-1)=x^2+x+1$ for all real numbers $x$. What is $f(5)$?


 $\mathrm {(A)} 7\qquad \mathrm {(B)} 13\qquad \mathrm {(C)} 31\qquad \mathrm {(D)} 111\qquad \mathrm {(E)} 211$

 \textbf{Answer: }A



\textbf{Solution}

$3x-1 =5 \implies x= 2$
$f(3(2)-1) = 2^2+2+1=7 \implies (A)$



	\item (\textit{AMC 12 12B 2007}) Some boys and girls are having a car wash to raise money for a class trip to China. Initially $40$\% of the group are girls. Shortly thereafter two girls leave and two boys arrive, and then $30$\% of the group are girls. How many girls were initially in the group?


 $\mathrm {(A)} 4\qquad \mathrm {(B)} 6\qquad \mathrm {(C)} 8\qquad \mathrm {(D)} 10\qquad \mathrm {(E)} 12$

 \textbf{Answer: }C



\textbf{Solution 1}

If we let $p$ be the number of people initially in the group, then $0.4p$ is the number of girls. If two girls leave and two boys arrive, the number of people in the group is still $p$, but the number of girls is $0.4p-2$. Since only $30\%$ of the group are girls,
 \begin{align*} \frac{0.4p-2}{p}&=\frac{3}{10}\\ 4p-20&=3p\\ p&=20\end{align*}
The number of girls initially in the group is $0.4p=0.4(20)=\boxed{\mathrm{(C) \ } 8}$




\textbf{Solution 2}

There are the same number of total people before and after, but the number of girls has dropped by two or $10\%$ of the total.  $\frac{2}{0.1}=20$, and $40\%\cdot20=8$, so the answer is $\mathrm{(C)}$.

 



\textbf{Solution 3}

Let $x$ be the number of people initially in the group and $g$ the number of girls. $\frac{2}{5}x = g$, so $x = \frac{5}{2}g$. Also, the problem states $\frac{3}{10}x = g-2$. Substituting $x$ in terms of $g$ into the second equation yields that $g = \boxed{ 8\ \mathrm{(C)}}$.

 


	\item (\textit{AMC 12 12B 2007}) The angles of quadrilateral $ABCD$ satisfy $\angle A=2\angle B=3\angle C=4\angle D$. What is the degree measure of $\angle A$, rounded to the nearest whole number?


 $\mathrm {(A)} 125\qquad \mathrm {(B)} 144\qquad \mathrm {(C)} 153\qquad \mathrm {(D)} 173\qquad \mathrm {(E)} 180$

 \textbf{Answer: }D



\textbf{Solution}

The sum of the interior angles of any quadrilateral is $360^\circ.$
 \begin{align*} 360 &= \angle A + \angle B + \angle C + \angle D\\ &= \angle A + \frac{1}{2}A + \frac{1}{3}A + \frac{1}{4}A\\ &= \frac{12}{12}A + \frac{6}{12}A + \frac{4}{12}A + \frac{3}{12}A\\ &= \frac{25}{12}A \end{align*}

 \[\angle A = 360 \cdot \frac{12}{25} = 172.8 \approx \boxed{\mathrm{(D) \ } 173}\]



	\item (\textit{AMC 12 12B 2007}) A teacher gave a test to a class in which $10\%$ of the students are juniors and $90\%$ are seniors. The average score on the test was $84$. The juniors all received the same score, and the average score of the seniors was $83$. What score did each of the juniors receive on the test?


 $\mathrm {(A)} 85\qquad \mathrm {(B)} 88\qquad \mathrm {(C)} 93\qquad \mathrm {(D)} 94\qquad \mathrm {(E)} 98$

 \textbf{Answer: }C



\textbf{Solution}

We can assume there are $10$ people in the class. Then there will be $1$ junior and $9$ seniors. The sum of everyone's scores is $10 \cdot 84 = 840$. Since the average score of the seniors was $83$, the sum of all the senior's scores is $9 \cdot 83 = 747$. The only score that has not been added to that is the junior's score, which is $840 - 747 = \boxed{\textbf{(C) } 93}$




\textbf{Solution 2}

Let the average score of the juniors be $j$. The problem states the average score of the seniors is $83$. The equation for the average score of the class (juniors and seniors combined) is $\frac{j}{10} + \frac{83 \cdot 9}{10} = 84$. Simplifying this equation yields $j = \boxed{\textbf{(C) } 93}$



	\item (\textit{AMC 12 12B 2007}) A traffic light runs repeatedly through the following cycle: green for $30$ seconds, then yellow for $3$ seconds, and then red for $30$ seconds. Leah picks a random three-second time interval to watch the light. What is the probability that the color changes while she is watching?


 $\mathrm {(A)} \frac{1}{63}\qquad \mathrm {(B)} \frac{1}{21}\qquad \mathrm {(C)} \frac{1}{10}\qquad \mathrm {(D)} \frac{1}{7}\qquad \mathrm {(E)} \frac{1}{3}$

 \textbf{Answer: }D



\textbf{Solution}

The traffic light runs through a $63$ second cycle.

 Letting $t=0$ reference the moment it turns green, the light changes at three different times: $t=30$, $t=33$, and $t=63$

 This means that the light will change if the beginning of Leah's interval lies in$[27,30]$, $[30,33]$ or $[60,63]$

 This gives a total of $9$ seconds out of $63$
$\frac{9}{63} = \frac{1}{7} \Rightarrow \mathrm{(D)}$



	\item (\textit{AMC 12 12B 2007}) Point $P$ is inside equilateral $\triangle ABC$. Points $Q$, $R$, and $S$ are the feet of the perpendiculars from $P$ to $\overline{AB}$, $\overline{BC}$, and $\overline{CA}$, respectively. Given that $PQ=1$, $PR=2$, and $PS=3$, what is $AB$?


 $\mathrm {(A)} 4\qquad \mathrm {(B)} 3\sqrt{3}\qquad \mathrm {(C)} 6\qquad \mathrm {(D)} 4\sqrt{3}\qquad \mathrm {(E)} 9$

 \textbf{Answer: }D



\textbf{Solution 1}

Drawing $\overline{PA}$, $\overline{PB}$, and $\overline{PC}$, $\triangle ABC$ is split into three smaller triangles. The altitudes of these triangles are given in the problem as $PQ$, $PR$, and $PS$.

 Summing the areas of each of these triangles and equating it to the area of the entire triangle, we get:

 
 \[\frac{s}{2} + \frac{2s}{2} + \frac{3s}{2} = \frac{s^2\sqrt{3}}{4}\]where $s$ is the length of a side of the equilateral triangle

 
 \[s = \boxed{\mathrm{(D) \ } 4\sqrt{3}}\]

            \begin{itemize}
        \item Note - This is called \href{https://artofproblemsolving.com/wiki/index.php?title=Viviani\%27s\_Theorem}{Viviani's Theorem}.
            \end{itemize}
            



\textbf{Solution 2}

The 60 degree angles suggest constructing 30-60-90 triangles. As such, let the foot of the altitude from $Q$ to $\overline{BC}$ be $D$. Also, since $\angle{PQD}=60$, let the foot of the altitude from $P$ to $\overline{QD}$ be $E$. Since $\triangle{QEP}$ is 30-60-90, and $PQ=1$, $QE=\frac{1}{2}$. Also, since $PR=2$, $DE=2$. Thus, $QD=\frac{5}{2}$. Once again, since $\triangle{QBD}$ is 30-60-90, $QB=\frac{5\sqrt{3}}{3}$. Similar reasoning to $\overline{AQ}$ and summing the segments yields $\boxed{(D) 4\sqrt{3}}$




\textbf{Solution 3 (Viviani's Theorem)}

According to Viviani's Theorem, in an equilateral triangle, the sum of the perpendicular distances from a point inside the triangle will equal the altitude of the triangle. Thus, the height of the triangle is 6. 

 Then we start to form a 30-60-90 triangle. By that, we divide $6/\sqrt{3} = 2\sqrt{3}$ then multiply 2, and we get an answer of $\boxed{(D) 4\sqrt{3}}$



	\item (\textit{AMC 12 12B 2007}) The geometric series $a+ar+ar^2\cdots$ has a sum of $7$, and the terms involving odd powers of $r$ have a sum of $3$. What is $a+r$?


 $\mathrm {(A)} \frac{4}{3}\qquad \mathrm {(B)} \frac{12}{7}\qquad \mathrm {(C)} \frac{3}{2}\qquad \mathrm {(D)} \frac{7}{3}\qquad \mathrm {(E)} \frac{5}{2}$

 \textbf{Answer: }E



\textbf{Solution 1}

The sum of an infinite geometric series is given by $\frac{a}{1-r}$ where $a$ is the first term and $r$ is the common ratio.

 In this series, $\frac{a}{1-r} = 7$

 The series with odd powers of $r$ is given as 
 \[ar + ar^3 + ar^5 ...\]

 Its sum can be given by $\frac{ar}{1-r^2} = 3$

 Doing a little algebra

 $ar = 3(1-r)(1+r)$
$ar = 3\left(\frac{a}{7}\right)(1+r)$
$\frac{7}{3}r = 1 + r$
$r = \frac{3}{4}$
$a = 7(1-r) = \frac{7}{4}$
$a + r =\boxed{ \frac{5}{2}} \Rightarrow \mathrm{(E)}$




\textbf{Solution 2}

The given series can be decomposed as follows: $(a + ar + ar^2 + \ldots) = (a + ar^2 + ar^4 + \ldots) + (ar + ar^3 + ar^5 + \ldots)$

 Clearly $(a + ar^2 + ar^4 + \ldots) = (ar + ar^3 + ar^5 + \ldots)/r = 3/r$. We obtain that $7 = 3/r + 3$, hence $r = \frac{3}{4}$.

 Then from $7 = (a + ar + ar^2 + \ldots) = \frac{a}{1-r}$ we get $a=\frac{7}{4}$, and thus $a + r = \boxed{\frac{5}{2}}$.

 


	\item (\textit{AMC 12 12B 2007}) Each face of a regular tetrahedron is painted either red, white, or blue. Two colorings are considered indistinguishable if two congruent tetrahedra with those colorings can be rotated so that their appearances are identical. How many distinguishable colorings are possible?


 $\mathrm {(A)} 15\qquad \mathrm {(B)} 18\qquad \mathrm {(C)} 27\qquad \mathrm {(D)} 54\qquad \mathrm {(E)} 81$

 \textbf{Answer: }A



\textbf{Solution 1}

A tetrahedron has 4 sides. We can ignore the rotation part of the question completely and focus on the colors.The ratio of the number of faces with each color must be one of the following:

 $4:0:0$, $3:1:0$, $2:2:0$, or $2:1:1$

 


 The first ratio yields $3$ appearances, one of each color.

 The second ratio yields $3\cdot 2 = 6$ appearances, three choices for the first color, and two choices for the second.

 The third ratio yields $\binom{3}{2} = 3$ appearances since the two colors are interchangeable.

 The fourth ratio yields $3$ appearances. There are three choices for the first color, and since the second two colors are interchangeable, there is only one distinguishable pair that fits them.

 The total is $3 + 6 + 3 + 3 = 15$ appearances $\Rightarrow \mathrm{(A)}$




\textbf{Solution 2}

Every colouring can be represented in the form $(w,r,b)$, where $w$ is the number of white faces, $r$ is the number of red faces, and $b$ is the number of blue faces. Every distinguishable colouring pattern can be represented like this in exactly one way, and every ordered whole number triple with a total sum of 4 represents exactly one colouring pattern (if two tetrahedra have rearranged colours on their faces, it is always possible to rotate one so that it matches the other).

 Therefore, the number of colourings is equal to the number of ways 3 distinguishable nonnegative integers can add to 4. If you have 6 cockroaches in a row, this number is equal to the number of ways to pick two of the cockroaches to eat for dinner (because the remaining cockroaches in between are separated in to three sections with a non-negative number of cockroaches each), which is $\binom{6}{2} = 15$

 Alternative explanation to solution 2:A regular tetrahedron is the only platonic solid in which any of the faces is adjacent to all the other 3 faces. Hence we only need to think about the number of faces we can colour for each face. Let the number of faces coloured with red, blue and white be $r, b, w$ respectively. So we are solving for the number of solutions of the equation:$r + b + w = 4$for nonnegative integers $r, b, w$. By Stars and Bars, we obtain the final answer which is $\binom{6}{2} = 15/(A)$



	\item (\textit{AMC 12 12B 2007}) If $a$ is a nonzero integer and $b$ is a positive number such that $ab^2=\log_{10}b$, what is the median of the set $\{0,1,a,b,1/b\}$?


 $\mathrm {(A)} 0\qquad \mathrm {(B)} 1\qquad \mathrm {(C)} a\qquad \mathrm {(D)} b\qquad \mathrm {(E)} \frac{1}{b}$

 \textbf{Answer: }D



\textbf{Solution}

Note that if $a$ is positive, then, the equation will have no solutions for $b$. This becomes more obvious by noting that at $a=1$, $ab^2 > \log_{10} b$. The LHS quadratic function will increase faster than the RHS logarithmic function, so they will never intersect.

 This puts $a$ as the smallest in the set since it must be negative.

 Checking the new equation: $-b^2 = \log_{10}b$

 Near $b=0$, $-b^2 > \log_{10} b$ but at $b=1$, $-b^2 < \log_{10} b$

 This implies that the solution occurs somewhere in between: $0 < b < 1$

 This also implies that $\frac{1}{b} > 1$

 This makes our set (ordered)$\{a,0,b,1,1/b\}$

 The median is $b \Rightarrow \mathrm {(D)}$




\textbf{Solution 2}

Let $b=0.1$. Then $a\cdot0.01 = -1,$ giving $a=-100$. Then the ordered set is $\{-100, 0, 0.1, 1, 10\}$ and the median is $0.1=b,$ so the answer is $\mathrm {(D)}$.

 


	\item (\textit{AMC 12 12B 2007}) Let $a$, $b$, and $c$ be digits with $a\ne 0$. The three-digit integer $abc$ lies one third of the way from the square of a positive integer to the square of the next larger integer. The integer $acb$ lies two thirds of the way between the same two squares. What is $a+b+c$?


 $\mathrm {(A)} 10\qquad \mathrm {(B)} 13\qquad \mathrm {(C)} 16\qquad \mathrm {(D)} 18\qquad \mathrm {(E)} 21$

 \textbf{Answer: }C



\textbf{Solution 1}

The difference between $acb$ and $abc$ is given by

 $(100a + 10c + b) - (100a + 10b + c) = 9(c-b)$

 The difference between the two squares is three times this amount or

 $27(c-b)$

 The difference between two consecutive squares is always an odd number, therefore $c-b$ is odd. We will show that $c-b$ must be 1. Otherwise we would be looking for two consecutive squares that are at least 81 apart. But already the equation $(x+1)^2-x^2 = 27\cdot 3$ solves to $x=40$, and $40^2$ has more than three digits.

 The consecutive squares with common difference $27$ are $13^2=169$ and $14^2=196$. One third of the way between them is $178$ and two thirds of the way is 

 


	\item (\textit{AMC 12 12B 2007}) Rhombus $ABCD$, with side length $6$, is rolled to form a cylinder of volume $6$ by taping $\overline{AB}$ to $\overline{DC}$. What is $\sin(\angle ABC)$?


 $\mathrm {(A)} \frac{\pi}{9}\qquad \mathrm {(B)} \frac{1}{2}\qquad \mathrm {(C)} \frac{\pi}{6}\qquad \mathrm {(D)} \frac{\pi}{4}\qquad \mathrm {(E)} \frac{\sqrt{3}}{2}$

 \textbf{Answer: }A



\textbf{Solution}


                    \begin{figure}[H]
                        \centering
                        \includegraphics[width=0.5\textwidth]{Solution/AMC_12_12B_2007_19_Solution_Figure_1.png}
                    \end{figure}
                    
$V_{\mathrm{Cylinder}} = \pi r^2 h$

 Where $C = 2\pi r = 6$ and $h=6\sin\theta$
$r = \frac{3}{\pi}$
$V = \pi \left(\frac{3}{\pi}\right)^2\cdot 6\sin\theta$
$6 = \frac{9}{\pi} \cdot 6\sin\theta$
$\sin\theta = \frac{\pi}{9} \Rightarrow \mathrm{(A)}$



	\item (\textit{AMC 12 12B 2007}) The parallelogram bounded by the lines $y=ax+c$, $y=ax+d$, $y=bx+c$, and $y=bx+d$ has area $18$. The parallelogram bounded by the lines $y=ax+c$, $y=ax-d$, $y=bx+c$, and $y=bx-d$ has area $72$. Given that $a$, $b$, $c$, and $d$ are positive integers, what is the smallest possible value of $a+b+c+d$?


 $\mathrm {(A)} 13\qquad \mathrm {(B)} 14\qquad \mathrm {(C)} 15\qquad \mathrm {(D)} 16\qquad \mathrm {(E)} 17$

 \textbf{Answer: }D



\textbf{Solution 1}

Plotting the parallelogram on the coordinate plane, the 4 corners are at $(0,c),(0,d),\left(\frac{d-c}{a-b},\frac{ad-bc}{a-b}\right),\left(\frac{c-d}{a-b},\frac{bc-ad}{a-b}\right)$. Because $72= 4\cdot 18$, we have that $4(c-d)\left(\frac{c-d}{a-b}\right) = (c+d)\left(\frac{c+d}{a-b}\right)$ or that $2(c-d)=c+d$, which gives $c=3d$ (consider a \href{https://artofproblemsolving.com/wiki/index.php?title=Homothety}{homothety}, or dilation, that carries the first parallelogram to the second parallelogram; because the area increases by $4\times$, it follows that the stretch along the diagonal, or the ratio of side lengths, is $2\times$). The area of the triangular half of the parallelogram on the right side of the y-axis is given by $9 = \frac{1}{2} (c-d)\left(\frac{d-c}{a-b}\right)$, so substituting $c = 3d$:

 
 \[\frac{1}{2} (c-d)\left(\frac{c-d}{a-b}\right) = 9 \quad \Longrightarrow \quad 2d^2 = 9(a-b)\]

 Thus $3|d$, and we verify that $d = 3$, $a-b = 2 \Longrightarrow a = 3, b = 1$ will give us a minimum value for $a+b+c+d$. Then $a+b+c+d = 3 + 1 + 9 + 3 = \boxed{\mathbf{(D)} 16}$.

 



\textbf{Solution 2}

The key to this solution is that area is invariant under translation. By suitably shifting the plane, the problem is mapped to the lines $c,d,(b-a)x+c,(b-a)x+d$ and $c,-d,(b-a)x+c,(b-a)x-d$. Now, the area of the parallelogram contained by is the former is equal to the area of a rectangle with sides $d-c$ and $\frac{d-c}{b-a}$, $\frac{(d-c)^2}{b-a}=18$, and the area contained by the latter is $\frac{(c+d)^2}{b-a}=72$. Thus, $d=3c$ and $b-a$ must be even if the former quantity is to equal $18$. $c^2=18(b-a)$ so $c$ is a multiple of $3$. Putting this all together, the minimal solution for $(a,b,c,d)=(3,1,3,9)$, so the sum is $\boxed{\textbf{(D)} 16}$.

 



\textbf{Solution 3}

Let $a$ and $b$ be the slopes of the lines such that $b > a$ (i.e. the line $bx+c$ is steeper than $ax+c$) and $c > d$ (i.e. the point $(0, c)$ is higher than the point $(0, d)$. Upon drawing a diagram, we see that both the smaller and the larger parallelogram can be split along the x-axis, such that both of their areas are the combinations of two corresponding triangles. The area of a triangle is $\frac{bh}{2}$, but since a parallelogram is two such triangles, the area becomes $bh$.

 Let $b_1$ and $h_1$ denote the base length and height, respectively, of the triangles pertaining to the smaller parallelogram and $b_2$ and $h_2$ denote those of the larger parallelogram. Notice that $b_1$ is simple the distance from $(0, d)$ to $(0, c)$, or $(c-d)$. Also notice that $h_1$ is the distance from the $x$-axis to the intersection of lines $ax+c$ and $bx+d$. This is equivalent to the value of the $x$-coordinate of intersection, so we solve for $x$:

 
$ax+c=bx+d$
$\Rightarrow bx-ax = c-d$
$\Rightarrow (b-a)x = c-d$
$\Rightarrow x = \frac{c-d}{b-a}$.

 
The area of the smaller parallelogram is $b_1*h_1$, or

 
$(c-d) * \frac{c-d}{b-a}$
$\Rightarrow \frac{(c-d)^2}{b-a}$.

 
$b_2$ is the distance from $(0, -d)$ to $(0, c)$, or $(c+d)$. $h_2$ is the $x$-coordinate of the intersection of the lines $ax+c$ and $bx-d$. Again, we solve for x:

 
$ax+c=bx-d$
$\Rightarrow bx-ax = c+d$
$\Rightarrow (b-a)x = c+d$
$\Rightarrow x = \frac{c+d}{b-a}$.

 
The area of the larger parallelogram is $b-1*h_1$, or

 
$(c+d) * \frac{c+d}{b-a}$
$\Rightarrow \frac{(c+d)^2}{b-a}$.

 
The areas of the parallelograms are given to us: $18$ and $72$. Therefore we can set up a ratio:

 
$\frac{18}{72} = \frac{\frac{(c-d)^2}{b-a}}{\frac{(c+d)^2}{b-a}}$
$\Rightarrow 18(c+d)^2 = 72(c-d)^2$
$\Rightarrow (c+d)^2 = 4(c-d)^2$
$\Rightarrow c^2 + 2cd + d^2 = 4c^2 - 8cd + 4d^2$
$\Rightarrow 3c^2 - 10cd +3d^2 = 0$
$\Rightarrow (3c-d)(c-3d)=0$
$\Rightarrow c=3d, c=\frac{d}{3}$

 
We established earlier that $c>d$, so $c=3d$. Plugging this into the intial equations yields

 
$\frac{16d^2}{b-a} = 72$

 
and

 
$\frac{4d^2}{b-a} = 18$

 Solving for $d$, we get

 
$d = 3\sqrt{\frac{b-a}{2}}$

 
We want the sum of $a$, $b$, $c$, and $d$. We can now rewrite this

 
$a + b + 12\sqrt{\frac{b-a}{2}}$

 
We are told that $a$, $b$, $c$, and $d$ are all positive integers. Therefore the value under the radical must be a perfect square greater than 0. We can rewrite this

 
$\frac{b-a}{2} = k^2$

 
where $k$ is some positive integer.

 Rearranging, we get

 
$b = a + 2k^2$

 Now we can rewrite the sum as

 
$a + a + 2k^2 +12k$.

 
Since both $a$ and $k$ must be at least $1$, the minimum value is

 
$1 + 1 + 2(1)^2 + 12(1) = 16 \Rightarrow \boxed{\text{D}}$.

 


	\item (\textit{AMC 12 12B 2007}) The first $2007$ positive integers are each written in base $3$. How many of these base-$3$ representations are palindromes? (A palindrome is a number that reads the same forward and backward.)


 $\mathrm {(A)} 100\qquad \mathrm {(B)} 101\qquad \mathrm {(C)} 102\qquad \mathrm {(D)} 103\qquad \mathrm {(E)} 104$

 \textbf{Answer: }A



\textbf{Solution 1}

$2007_{10} = 2202100_{3}$

 All numbers of six or less digits in base 3 have been written.

 The form of each palindrome is as follows

 1 digit  - $a$

 2 digits - $aa$

 3 digits - $aba$

 4 digits - $abba$

 5 digits - $abcba$

 6 digits - $abccba$

 Where $a,b,c$ are base 3 digits

 Since $a \neq 0$, this gives a total of $2 + 2 + 2\cdot 3 + 2\cdot 3 + 2\cdot 3^2 + 2\cdot 3^2 = 52$ palindromes so far.

 7 digits - $abcdcba$, but not all of the numbers are less than $2202100_{3}$

 Case: $a=1$

 All of these numbers are less than $2202100_3$ giving $3^3$ more palindromes

 Case: $a=2$, $b\neq 2$

 All of these numbers are also small enough, giving $2\cdot 3^2$ more palindromes

 Case: $a=2$, $b=2$

 It follows that $c=0$, since any other $c$ would make the value too large. This leaves the number as $220d022_3$. Checking each value of d, all of the three are small enough, so that gives $3$ more palindromes.

 Summing our cases there are$52 + 3^3 + 2\cdot 3^2 + 3 = 100 \Rightarrow \mathrm{(A)}$




\textbf{Solution 2 (similar to 1, no cases)}

Notice 2017 (base 3) is 2202100. This means we aim to find palindromes up to 2202100 using only digits (0,1,2) as we are in base 3. Therefore, there are 3 options for each letter (with the exception of a, which has 2 options; this is because we don't want to consider leading zeroes as this only complicates things).Then, consider:1 digit palindromes are of the form $a$,2 digit of the form $aa$,3 of the form $aba$,4 of the form $abba$,and so on.

 Since some digits repeat, it suffices to call them of the form$a$, $a$,$ab$, $ab$,$abc$, $abc$,for palindromes up to six digits. This gives us $(2, 2, 6, 6, 18, 18)$ options respectively, or 106 total.

 For seven digits, we aim to find all 7-digit palindromes, minus the ones greater than 2202100.There are $2\cdot3\cdot3\cdot3$ options for 7-digit palindromes (as they are of the form $abcd$), and if we aim to find ones greater than 2202100, we know the last two digits are both 2 (considering there are no options greater than 22 for the first two). Therefore we are left with $aba$, with 2 options for $a$ (because there are no palindromes above 021 with $a = 0$), and 3 options for $b$, giving 6 options total.Therefore the total number of base-3 palindromes below 2202100 is $106-6=\boxed{\mathrm{(A)}{100}}$



	\item (\textit{AMC 12 12B 2007}) Two particles move along the edges of equilateral $\triangle ABC$ in the direction 
 \[A\Rightarrow B\Rightarrow C\Rightarrow A,\] starting simultaneously and moving at the same speed. One starts at $A$, and the other starts at the midpoint of $\overline{BC}$. The midpoint of the line segment joining the two particles traces out a path that encloses a region $R$. What is the ratio of the area of $R$ to the area of $\triangle ABC$?


 $\mathrm {(A)} \frac{1}{16}\qquad \mathrm {(B)} \frac{1}{12}\qquad \mathrm {(C)} \frac{1}{9}\qquad \mathrm {(D)} \frac{1}{6}\qquad \mathrm {(E)} \frac{1}{4}$

 \textbf{Answer: }A



\textbf{Solution 1}

First, notice that each of the midpoints of $AB$,$BC$, and $CA$ are on the locus.Suppose after some time the particles have each been displaced by a short distance $x$, to new positions $A'$ and $M'$ respectively. Consider $\triangle ABM$ and drop a perpendicular from $M'$ to hit $AB$ at $Y$. Then, $BA'=1-x$ and $BM'=1/2+x$. From here, we can use properties of a $30-60-90$ triangle to determine the lengths $YA'$ and $YM'$ as monomials in $x$. Thus, the locus of the midpoint will be linear between each of the three special points mentioned above. It follows that the locus consists of the only triangle with those three points as vertices. (A cheaper way to find the shape of the region is to look at the answer choices: if it were any sort of conic section then the ratio would not generally be rational.)Comparing inradii between this "midpoint" triangle and the original triangle, the area contained by $R$ must be $\textbf{(A)} \frac{1}{16}$ of the total area.

 


 



\textbf{Solution 2}

\href{https://artofproblemsolving.com/wiki/index.php?title=WLOG}{WLOG}, let point $A$ be located on the origin. Let side $\overline{AB}$ be on the $x$-axis, and let point $B$ be $(0, 1)$.When one of the ants is on a vertex, the other will be on the midpoint of the opposite side. Marking these points gives us a rough view of what seem to be the vertices of a smaller equilateral triangle, that is concentric with the larger one. After observing(or imagining) the ants moving, we find that the area enclosed is indeed an equilateral triangle.

 Let the midpoint of $\overline{BC}$ be $D$. Then we can find the coordinates of $D$ using the equations of two lines: the extension of $\overline{AD}$ and the extension of $\overline{BC}$. The slope of any line is simply the tangent of the angle between the line and the $x$-axis. The angle between $\overline{AD}$ and the $x$-axis is $30^{\circ}$, and $tan(30) = \frac{\sqrt{3}}{3}$.

 Therefore the equation of the line containing $\overline{AD}$ is $y= \frac{\sqrt{3}}{3}x.$

 The slope of the line containing $\overline{BC}$ is $-\sqrt{3}$, and it goes through the point $(1, 0)$. After solving for the $y$-intercept, we get $y = -\sqrt{3}x+\sqrt{3}$. Setting these equal to each other yields

 $D\left(\frac{3}{4}, \frac{\sqrt{3}}{4}\right).$

 The midpoint of $D$ and the $A$ is the bottom-left vertex of the smaller equilateral triangle. Using the \href{https://artofproblemsolving.com/wiki/index.php?title=Distance\_formula}{distance formula}, we get the point

 $\left(\frac{3}{8}, \frac{\sqrt{3}}{8}\right).$

 Since the smaller triangle is simply a dilation of the larger one, they both have an axis of symmetry over $x=\frac{1}{2}$. This yields the bottom-right vertex of the smaller triangle as

 $\left(\frac{5}{8}, \frac{\sqrt{3}}{8}\right).$

 Therefore the side length of the smaller triangle is $\frac{1}{4}.$ The problem asks us for the ratio of the areas, and since area is two-dimensional, we simply square the ratio to get $\frac{1}{16} \Rightarrow \boxed{\text{A}}.$



	\item (\textit{AMC 12 12B 2007}) How many non-congruent right triangles with positive integer leg lengths have areas that are numerically equal to $3$ times their perimeters?


 $\mathrm {(A)} 6\qquad \mathrm {(B)} 7\qquad \mathrm {(C)} 8\qquad \mathrm {(D)} 10\qquad \mathrm {(E)} 12$

 \textbf{Answer: }A



\textbf{Solution 1}

Let $a$ and $b$ be the two legs of the triangle.

 We have $\frac{1}{2}ab = 3(a+b+c)$.

 Then $ab=6 \left(a+b+\sqrt {a^2 + b^2}\right)$.

 We can complete the square under the root, and we get, $ab=6 \left(a+b+\sqrt {(a+b)^2 - 2ab}\right)$.

 Let $ab=p$ and $a+b=s$, we have $p=6 \left(s+ \sqrt {s^2 - 2p}\right)$.

 After rearranging, squaring both sides, and simplifying, we have $p=12s-72$.

 
Putting back $a$ and $b$, and after factoring using Simon's Favorite Factoring Trick, we've got $(a-12)(b-12)=72$.

 
Factoring 72, we get 6 pairs of $a$ and $b$
$(13, 84), (14, 48), (15, 36), (16, 30), (18, 24), (20, 21).$

 
And this gives us $6$ solutions $\Rightarrow \mathrm{(A)}$.

 
Alternatively, note that $72 = 2^3 \cdot 3^2$. Then 72 has $(3+1)(2+1) = (4)(3) = 12$ factors. However, half of these are repeats, so we have $\frac{12}{2} = 6$ solutions.

 



\textbf{Solution 2}

We will proceed by using the fact that $[ABC] = r\cdot s$, where $r$ is the radius of the incircle and $s$ is the semiperimeter $\left(s = \frac{p}{2}\right)$.

 We are given $[ABC] = 3p = 6s \Rightarrow rs = 6s \Rightarrow r = 6$.

 The incircle of $ABC$ breaks the triangle's sides into segments such that $AB = x + y$, $BC = x + z$ and $AC = y + z$. Since ABC is a right triangle, one of $x$, $y$ and $z$ is equal to its radius, 6. Let's assume $z = 6$.

 The side lengths then become $AB = x + y$, $BC = x + 6$ and $AC = y + 6$. Plugging into Pythagorean's theorem:

 $(x + y)^2 = (x+6)^2 + (y + 6)^2$
$x^2 + 2xy + y^2 = x^2 + 12x + 36 + y^2 + 12y + 36$
$2xy - 12x - 12y = 72$
$xy - 6x - 6y = 36$
$(x - 6)(y - 6) - 36 = 36$
$(x - 6)(y - 6) = 72$

 We can factor $72$ to arrive with $6$ pairs of solutions: $(7, 78), (8,42), (9, 30), (10, 24), (12, 18),$ and $(14, 15) \Rightarrow \mathrm{(A)}$.

 



\textbf{Solution 3}

Let $a$ and $b$ be the two legs of the triangle, and $c$ be the hypotenuse.

 By using $Area = \frac{r}{2} (a+b+c)$, where $r$ is the in-radius, we get:

 
 \[3(a+b+c) = \frac{r}{2} (a+b+c)\]
 \[r=6\]

 In right triangle, $r = \frac{a+b-c}{2}$
 \[a+b-c = 12\]
 \[c = a + b - 12\]

 
By the triangle's area we get:

 
 \[\frac{ab}{2} = 6 \cdot \frac{a+b+c}{2}\]
 \[ab = 6(a+b+c)\]

 By substituting $c$ in:

 
 \[ab = 6(a+b+a + b - 12)\]
 \[ab - 12a - 12b + 72 = 0\]
 \[(a - 12)(b - 12) = 72\]

 As $72 = 2^3 \cdot 3^2$, there are $\frac{(3+1)(2+1)}{2} = 6$ solutions, $\boxed{\textbf{(A) } 6}$.

 



\textbf{Solution 4}

All pythagorean triples can be parametrized in the form $(a, b, c) = k(r^2 - s^2), k(2rs), k(r^2 + s^2)$ for positive integers $k, r, s$. The area being triple the perimeter implies that 
 \[k^2(r^2 - s^2)rs = 3(k(r^2 - s^2) + k(2rs) + k(r^2 + s^2)).\] This can be simplified to get 
 \[ks(r - s) = 6.\] Now, we get the triples 
 \[(k, r, s) = (1, 7, 1), (1, 5, 2), (1, 5, 3), (1, 7, 6), (2, 4, 1), (2, 4, 3), (3, 3, 1), (3, 3, 2), (6, 2, 1).\] However, the ones where $r$ and $s$ are not different signs and relatively prime are redundant, so we get $6$ triples total.

 


	\item (\textit{AMC 12 12B 2007}) How many pairs of positive integers $(a,b)$ are there such that $\gcd{(a,b)}=1$ and 
 \[\frac{a}{b}+\frac{14b}{9a}\] is an integer?


 $\mathrm {(A)} 4\qquad \mathrm {(B)} 6\qquad \mathrm {(C)} 9\qquad \mathrm {(D)} 12\qquad \mathrm {(E)} \text{infinitely many}$

 \textbf{Answer: }A



\textbf{Solution 1}

Combining the fraction, $\frac{9a^2 + 14b^2}{9ab}$ must be an integer.

 Since the denominator contains a factor of $9$, $9 | 9a^2 + 14b^2 \quad\Longrightarrow\quad 9 | b^2 \quad\Longrightarrow\quad 3 | b$

 Since $b = 3n$ for some positive integer $n$, we can rewrite the fraction(divide by $9$ on both top and bottom) as $\frac{a^2 + 14n^2}{3an}$

 Since the denominator now contains a factor of $n$, we get $n | a^2 + 14n^2 \quad\Longrightarrow\quad n | a^2$.

 But since $1=\gcd(a,b)=\gcd(a,3n)=\gcd(a,n)$, we must have $n=1$, and thus $b=3$.

 For $b=3$ the original fraction simplifies to $\frac{a^2 + 14}{3a}$.

 For that to be an integer, $a$ must be a factor of $14$, and therefore we must have $a\in\{1,2,7,14\}$. Each of these values does indeed yield an integer.

 Thus there are four solutions: $(1,3)$, $(2,3)$, $(7,3)$, $(14,3)$ and the answer is $\mathrm{(A)}$




\textbf{Solution 2}

Let's assume that $\frac{a}{b} + \frac{14b}{9a} = m$ We get

 $9a^2 - 9mab + 14b^2 = 0$

 Factoring this, we get $4$ equations-

 $(3a-2b)(3a-7b) = 0$
$(3a-b)(3a-14b) = 0$
$(a-2b)(9a-7b) = 0$
$(a-b)(9a-14b) = 0$

 (It's all negative, because if we had positive signs, $a$ would be the opposite sign of $b$)

 Now we look at these, and see that-

 $3a=2b$
$3a=b$
$3a=7b$
$3a=14b$
$a=2b$
$9a=7b$
$a=b$
$9a=14b$

 This gives us $8$ solutions, but we note that the middle term needs to give you back $9m$.

 For example, in the case 

 $(a-2b)(9a-7b)$, the middle term is $-25ab$, which is not equal by $-9m$ for any integer $m$.

 Similar reason for the fourth equation. This eliminates the last four solutions out of the above eight listed, giving us 4 solutions total $\mathrm {(A)}$




\textbf{Solution 3}

Let $u = \frac{a}{b}$. Then the given equation becomes $u + \frac{14}{9u} = \frac{9u^2 + 14}{9u}$.

 Let's set this equal to some value, $k \Rightarrow \frac{9u^2 + 14}{9u} = k$.

 Clearing the denominator and simplifying, we get a quadratic in terms of $u$:

 $9u^2 - 9ku + 14 = 0 \Rightarrow u = \frac{9k \pm \sqrt{(9k)^2 - 504}}{18}$

 Since $a$ and $b$ are integers, $u$ is a rational number. This means that $\sqrt{(9k)^2 - 504}$ is an integer.

 Let $\sqrt{(9k)^2 - 504} = x$. Squaring and rearranging yields:

 $(9k)^2 - x^2 = 504$
$(9k+x)(9k-x) = 504$.

 In order for both $x$ and $a$ to be an integer, $9k + x$ and $9k - x$ must both be odd or even. (This is easily proven using modular arithmetic.) In the case of this problem, both must be even. Let $9k + x = 2m$ and $9k - x = 2n$.

 Then:

 $2m \cdot 2n = 504$
$mn = 126$.

 Factoring 126, we get $6$ pairs of numbers: $(1,126), (2,63), (3,42), (6,21), (7,18),$ and $(9,14)$.

 Looking back at our equations for $m$ and $n$, we can solve for $k = \frac{2m + 2n}{18} = \frac{m+n}{9}$. Since $k$ is an integer, there are only $2$ pairs of $(m,n)$ that work: $(3,42)$ and $(6,21)$. This means that there are $2$ values of $k$ such that $u$ is an integer. But looking back at $u$ in terms of $k$, we have $\pm$, meaning that there are $2$ values of $u$ for every $k$. Thus, the answer is $2 \cdot 2 = 4 \Rightarrow \mathrm{(A)}$.

 



\textbf{Solution 4}

Rewriting the expression over a common denominator yields $\frac{9a^2 + 14b^2}{9ab}$. This expression must be equal to some integer $m$.

 Thus, $\frac{9a^2 + 14b^2}{9ab} = m \rightarrow 9a^2 + 14b^2 = 9abm$. Taking this $\pmod{a}$ yields $14b^2 \equiv 0\pmod{a}$. Since $\gcd(a,b)=1$, $14 \equiv 0\pmod{a}$. This implies that $a|14$ so $a = 1, 2, 7, 14$.

 We can then take $9a^2 + 14b^2 = 9abm \pmod{b}$ to get that $9 \equiv 0 \pmod{b}$. Thus $b = 1, 3, 9$.

 However, taking $9a^2 + 14b^2 = 9abm \pmod{3}$, $b^2 \equiv 0\pmod{3}$ so $b$ cannot equal 1.

 Also, note that if $b = 9$,  $\frac{a}{b}+\frac{14b}{9a} = \frac{a}{9}+\frac{14}{a}$. Since $a|14$, $\frac{14}{a}$ will be an integer, but $\frac{a}{9}$ will not be an integer since none of the possible values of $a$ are multiples of 9. Thus, $b$ cannot equal 9.

 Thus, the only possible values of $b$ is 3, and $a$ can be 1, 2, 7, or 14. This yields 4 possible solutions, so the answer is $\mathrm{(A)}$.

 



\textbf{Solution 5 (Similar to Solution 1)}

Rewriting $\frac{a}{b} + \frac{14b}{9a}$ over a common denominator gives $\frac{9a^2 + 14b^2}{9ab}.$

 Thus, we have $9 \mid 9a^2 + 14b^2 \Rightarrow  3 \mid b.$

 Next, we have $ab \mid 9a^2+14b^2 \Rightarrow  ab \mid 14b^2 \Rightarrow a \mid 14b.$

 Thus, $a \in (1,2,7,4).$ 

 Next, we have $b \mid 9a^2 + 14b^2 \Rightarrow b \mid 9a^2 \Rightarrow b \mid 9.$

 Thus, $b \in (1,3,9).$

 Now, we simply do casework on $b.$ 

 Plugging in $b = 1,3$ and $9$ gives that there are $4$ total solutions for $(a,b).$




\textbf{Solution 6 (Similar to solution 3)}

Let $\frac{a}{b} = r.$ So $r + \frac{14}{9r}=I$, where I is an integer. Algebraic manipulations yield: $r^2-Ir+\frac{14}{9}=0$. The discriminant of this must be the square of a rational number, call this R. So $I^2-\frac{56}{9}=R^2 \longrightarrow I^2-R^2=(I-R)(I+R)=\frac{56}{9}$. I is $\frac{1}{2}$ the sum of $I-R$ and $I+R$. To have an integer sum, $I-R$ and $I+R$ must have the same denominator, namely 3. We proceed with casework.

 Case 1. $I+R=56/3$, $I-R=1/3$. This yields $I=19/2$, which is not an integer. This case produces 0 solutions.

 Case 2.$I+R=28/3$, $I-R=2/3$. This yields $I=5$. Substituting into our original equation yields: $r^2-5r+\frac{14}{9}=0$. Factoring gives: $r=\frac{1}{3}$, $r=\frac{14}{3}$. This case produces 2 solutions, namely (1,3) and (14,3).

 Case 3. $I+R=14/3$, $I-R=4/3$. This yields $I=3$. Substituting into our original equation yields: $r^2-3r+\frac{14}{9}=0$. Factoring gives: $r=\frac{2}{3}$, $r=\frac{7}{3}$. This case produces 2 solutions, namely (2,3) and (7,3).

 Case 4. $I+R=8/3$, $I-R=7/3$. This yields $I=5/2$, which is not an integer. This case produces 0 solutions.

 Altogether, we have 4 solutions, so our answer is $\boxed{(A)}$.

 



\textbf{Solution 7}

Rewrite the equation$\frac{a}{b}+\frac{14b}{9a}=k$in two different forms. First, multiply both sides by $b$ and subtract $a$ to obtain$\frac{14b^2}{9a}=bk-a.$Because $a$, $b$, and $k$ are integers, $14b^2$ must be a multiple of $a$, and because $a$ and $b$ have no common factors greater than 1, it follows that 14 is divisible by $a$. Next, multiply both sides of the original equation by $9a$ and subtract $14b$ to obtain$\frac{9a^2}{b}=9ak-14b.$This shows that $9a^2$ is a multiple of $b$, so 9 must be divisible by $b$. Thus if $(a,b)$ is a solution, then $b=1$, $3$, or $9$, and $a=1$, 2, 7, or 14. This gives a total of twelve possible solutions $(a,b)$, each of which can be checked quickly.$\frac{a}{b}+\frac{14b}{9a}$will only be an integer when $(a,b) \in \{(1,3),(2,3), (7,3), (14,3)\},$ for a total of $\boxed{4}$ pairs.

 


	\item (\textit{AMC 12 12B 2007}) Points $A,B,C,D$ and $E$ are located in 3-dimensional space with $AB=BC=CD=DE=EA=2$ and $\angle ABC=\angle CDE=\angle DEA=90^o$. The plane of $\triangle ABC$ is parallel to $\overline{DE}$. What is the area of $\triangle BDE$?


 $\mathrm {(A)} \sqrt{2}\qquad \mathrm {(B)} \sqrt{3}\qquad \mathrm {(C)} 2\qquad \mathrm {(D)} \sqrt{5}\qquad \mathrm {(E)} \sqrt{6}$

 \textbf{Answer: }C



\textbf{Solution 1}


                    \begin{figure}[H]
                        \centering
                        \includegraphics[width=0.5\textwidth]{Solution/AMC_12_12B_2007_25_Solution_Figure_1.png}
                    \end{figure}
                    Link to graph: \href{https://www.math3d.org/pHFSD6vRi}{https://www.math3d.org/pHFSD6vRi}

 
Let $A=(0,0,0)$, and $B=(2,0,0)$. Since $EA=2$, we could let $C=(2,0,2)$, $D=(2,2,2)$, and $E=(2,2,0)$. Now to get back to $A$ we need another vertex $F=(0,2,0)$. Now if we look at this configuration as if it was two dimensions, we would see a square missing a side if we don't draw $FA$. Now we can bend these three sides into an equilateral triangle, and the coordinates change: $A=(0,0,0)$, $B=(2,0,0)$, $C=(2,0,2)$, $D=(1,\sqrt{3},2)$, and $E=(1,\sqrt{3},0)$. Checking for all the requirements, they are all satisfied. Now we find the area of triangle $BDE$. The side lengths of this triangle are $2, 2, 2\sqrt{2}$, which is an isosceles right triangle. Thus the area of it is $\frac{2\cdot2}{2}=2\Rightarrow \mathrm{(C)}$.

 



\textbf{Solution 2}

Similar to solution 1, we allow $A=(0,0,0)$, $B=(2,0,0)$, and $C=(0,2,0)$. This creates the isosceles right triangle on the plane of $z=0$

 Now, note that $\angle CDE=\angle DEA=90^o$. This means that there exists some vector $DE$ parallel to the plane of $ABC$ that forms two right angles with $AE$ and $CD$. By definition, this is the cross product of the two vectors $AE$ and $CD$. Finding this cross product, we take the determinant of vectors

 $AE=<x_1,y_1,z>$ and

 $CD=<x_2,y_2,z>$ *Note that z is constant because the line is parallel to the plane*

 to get $(y_1-y_2)zi+(x_1-x_2)zj+(x_1y_2-y_1x_2)k$

 Because there can be no movement in the $z$ direction, the k unit vector must be zero. Also, because the i unit vector must be orthogonal and also 0. Thus, the vector of line $DE$ is simply $2tj+2k$

 From this, you can figure out that line $BE=2$, and the area of $BDE=\frac{2\cdot2}{2}=2\Rightarrow \mathrm{(C)}$.

 


\end{enumerate}